\begin{remark}[Arity at the zero-sum grade; AI-checked, pending human review]
\label{rem:fc-zero-sum-arity}
This remark comes from an AI-only pass (\file{research/audits/reviews/2026-11-05-rees-s3-arity.md}, with code in its folder), checked by a second, adversarial AI pass recorded in the same report; no human has checked it, nothing else in the paper uses it, and (R1) is unchanged. Keep the notation of Remark~\ref{rem:fc-zero-sum-sis}. Every statement concerns banks whose labels are independent and uniform on $K_d^\perp$, as in Lemma~\ref{lem:fc-forest}(a), and binary weights, since $s^d \in \{0,1\}$ at $d = p-1$.

\textup{(a)} \emph{A node of polynomial arity with degree $1+\lceil\log_2c\rceil$.} Put $j^* = \lceil\log_2(p-1)\rceil$, $D^{(0)} = \{0,1\}$ and $D^{(j)} = \{-2^{j-1},\dots,2^{j-1}\}$ for $j \ge 1$, read modulo $p$. Then $D^{(j^*)}$ contains every residue, and every $\alpha \in D^{(j+1)}$ equals $x - x'$ with $x = \lceil\alpha/2\rceil$ and $x' = -\lfloor\alpha/2\rfloor$ in $D^{(j)}$. A \emph{level-$j$ solver} in dimension $n$ with target $t \in \Fp^n$ reads independent uniform columns $v_i \in \Fp^n$ one at a time and returns a nonzero weight vector with entries in $D^{(j)}$ on the columns it read, with weighted sum $t$. For $j = j^*$ (used with $t = 0$ only) it returns a kernel vector of $n+1$ columns. For $n = 0$ it returns one column with weight $1$. For $j < j^*$ and $n \ge 1$ it repeats the following until the subset sums of $\beta_1,\dots,\beta_K$ cover $\Fp$: run the level-$(j+1)$ solver with target $0$ on the first $n-1$ coordinates of fresh columns, obtaining $\alpha_k$, and put $\beta_k = \sum_i\alpha_{k,i}v_{i,n}$. It then calls itself in dimension $n-1$, on fresh columns, with target the first $n-1$ coordinates of $t - c$, where $c = \sum_k\sum_i x'_{k,i}v_i$. Finally, on a set of indices $k$ whose $\beta_k$ sum to the remaining last coordinate, it uses $x_k$ in place of $x'_k$. The inner solver reads only the first $n-1$ coordinates of its columns, so given everything else each $\beta_k$ is uniform. Hence the solver halts with probability~$1$, its output is checked exactly, and by Wald's identity its expected number of columns satisfies $N_{j^*}(n) = n+1$ and $N_j(n) = 1 + s_p\sum_{e<n}N_{j+1}(e)$. Here $s_p \le \min(p,\lceil2\log_2(p-1)\rceil+4)$ is the expected number of uniform elements of $\Fp$ drawn until their subset sums cover $\Fp$. So $N_0(n) \le (s_p+1)^{j^*}\binom{n+j^*+1}{j^*+1}$, and the time is $O(R(n^2+p))$ operations in $\Fp$ when $R$ columns are read.

A node at $d \in \Gamma_0$ spans $W$ by its first $g_d + \lceil\log_p(2/\tau)\rceil$ labels and writes later labels in an echelon basis of $W$. It then runs the level-$0$ solver on up to $\lceil\log_2(2/\tau)\rceil$ disjoint blocks of $\lceil2N_0(g_d)\rceil$ labels. It checks as in (F5) and converts as in (F6). By Lemma~\ref{lem:s3-moment}(d) and Markov's inequality it aborts with probability at most $\tau$. With $\tau = \tau_d$ as in the variant after Lemma~\ref{lem:fc-zero-sum}, the bound $\eps/9$ of Lemma~\ref{lem:fc-forest}(b) is unchanged. Since $c = p-1$, the arity is $O_c(g_d^{\,1+\lceil\log_2c\rceil}\log(1/\tau_d))$. For comparison, $B_c(g_d)$ has degree $c\lceil\log_2(c+1)\rceil$; for example, the degrees are $3$ and $12$ at $c = 4$. For $p \ge 5$, the same step applied once, on top of a worst-case signed solver, is~\cite[Theorem~1.10]{KothariMetgerODonnellShuttyWu2026}, of degree $1+\lfloor(c+4)/4\rfloor$ with a factor $p^{p^2}$ and failure probability $2^{-\Omega(g_d)}$. At constant $\tau$ the degree here is smaller for $p = c+1 = 17$ and every prime $p \ge 23$, equal for $p \in \{3,5,13,19\}$, and larger for $p \in \{7,11\}$. At $\tau = 2^{-g_d}$ the factor $\log(1/\tau)$ adds one to the degree, which is then smaller only for primes $p \ge 29$.

The node is not total: if one coordinate of every label in the echelon basis repeats another, every $\beta_k$ vanishes and the node aborts. It gives neither arity $o(g_d^2)$ (at $p = 3$ it uses about $1.5g_d^2$ labels) nor time and arity $f(p)\poly(g_d)$. In the regime $p = O(\log m)$ of Remark~\ref{rem:fc-zero-sum-sis}(d) it is polynomial when $\log g_d\cdot\log p = O(\log m)$. That range is wider than the range $g_d\log p = O(\log m)$ of Lemma~\ref{lem:fc-zero-sum}(c), but it does not cover every profile.

\textup{(b)} \emph{Why the natural routes to arity $o(g_d^2)$ fail.} Write $n = a_d$, $\gamma = 2p/(p+1)$, and let $\delta$ be the success probability. These statements hold whatever the running time.
\begin{enumerate}[label=(\roman*),leftmargin=2em]
  \item \emph{Merging.} Consider an algorithm that works in stages $j = 1,\dots,\ell$. In stage $j$ it chooses, from what it has read, a linear map $\varphi_j$ that raises the rank of $(\varphi_1,\dots,\varphi_j)$ by $b_j \ge 1$, and it reads $\varphi_j$ of every label. It then replaces its family of pairwise disjoint sets of labels (initially the singletons) by pairwise disjoint unions of members whose $\varphi_j$-sum vanishes. BKW-type merging and disjoint-matching variants of Wagner's merging with binary weights, and the block recursion with Davenport groups, are of this form. Wagner's $k$-tree itself reuses elements across matches and is not. Given the past, the $\varphi_j$-sums of the members are independent and uniform. So the expected number of members drops by a factor of at least $\gamma$ per stage, and success with at least $L$ stages has probability at most $M\gamma^{-L}$. If every union has at most $F$ members, a stage with $b_j \ge b$ is nonempty with probability at most $(M+1)^Fp^{-b}$. Hence $(F\ln(M+1)+\ln(2n/\delta))\ln(2M/\delta) \ge n\ln p\ln\gamma$, that is, $M \ge \exp(\Omega(\sqrt{n\ln p\ln\gamma/F}))$. With unbounded $F$ and polynomial $M$, some stage must solve the same average-case problem in dimension at least $n\ln\gamma/\ln(2M/\delta)$.
  \item \emph{Switch recursion.} Consider schemes of the form of~(a) at level $0$, along a coordinate flag, with switches found by any rule, including sparse or exponential-time rules. All switches are fixed before the backward pass, and every action at dimension $d$ depends only on coordinates at most $d$ of the labels. Given the switches, each residual coordinate is uniform. A nonzero $\{0,\pm1\}$ relation on at most $((d-1)\ln p - \ln(4n/\delta))/\ln(2M)$ of $M$ uniform vectors of $\Fp^{d-1}$ exists with small probability. Hence such a scheme that reads $M$ labels satisfies $M\ln(2M) \ge \lceil\log_2p\rceil\,(\ln p\,L(L-1)/2 - L\ln(4n/\delta))$, where $L = n - \lfloor\log_2(4/\delta)/(\log_2p-\lfloor\log_2p\rfloor)\rfloor$. So $M = \Omega_p(n^2/\log n)$.
  \item \emph{Oblivious information sets.} Fix an information set and try $T$ binary patterns on its free labels, chosen without reading the free labels and not adapted to earlier trials. This succeeds with probability at most $T(2/p)^n$, the analogue of Proposition~\ref{prop:rc-lower}.
  \item \emph{Zero-sum theorems.} The Davenport bound of Lemma~\ref{lem:fc-zero-sum}(b) and its Chevalley--Warning proof give totality and membership in $\mathrm{PPA}_p$~\cite{GoosKamathSotirakiZampetakis2020}, but no search procedure. In the average case a zero-sum exists already at about $n\log_2p$ labels.
\end{enumerate}
None of this excludes reusing labels across dependencies as in~\cite{KothariODonnellWu2026}, which reaches about $a^2/3$ labels at $p = 3$. Whether average-case $\mathrm{BinSIS}$ at fixed $p$ is solvable in polynomial time at arity $o(a_d^2)$, or in time and arity $f(p)\poly(a_d)$, remains open.
% NEEDS HUMAN VERIFICATION: Theorem 1.10 of KothariMetgerODonnellShuttyWu2026 (arXiv 2609.40321) and its exponent 1+floor((q+3)/4); novelty of (a); no bib entries exist for Blum--Kalai--Wasserman or Wagner, so none are cited.
\end{remark}
